% metrics.tex -- Methodology: metrics (draft 4 Oct 2026). The definitions are the ones used by the
% analysis tools (multi_drone_control/tools/method_metrics.py, push_metrics.py).
% 9 Oct (structure audit): the unused "carried fraction" metric cut; "every Gazebo case is flown at
% least twice" replaced, since several cells are single runs (marked (1)).

\subsubsection{Metrics}
\label{sec:metrics}

All metrics are computed from the logged motion-capture poses of the payload and the UAVs, with the
events (TAKEOFF, the push, DETACH, LAND) from the run's event log.
\begin{description}
    \item[Height error] $e_z = z_L - z_L^{\mathrm{ref}}$, positive with the payload high; in hover, averaged
    over the 15\,s before LAND. \textbf{Heave:} its peak-to-peak over the hover window.
    \item[Horizontal error] $\|\mathbf{e}_{xy}\|$, the horizontal distance between the payload's centre and
    its reference; on a trajectory, its mean over the first lap or the whole figure, and its maximum.
    \item[Path error and lag] On a trajectory the error is split three ways, because a single number mixes
    them. The \textbf{ring RMSE} frozen for the carry claim is the ring's distance from the planner's reference as
    published once per 10\,Hz tick and held between ticks, from 5\,s after TAKEOFF to LAND (the \textbf{path RMSE} is
    the same over the path alone); the held reference is a
    staircase behind the schedule, so a ring exactly on schedule pays $v \cdot 0.1/\sqrt{3}$, 17\,mm at
    0.3\,m/s. The \textbf{clock-free error} is the ring's signed distance to the path it was sent, positive inside
    a curve: its mean over the path's tightest quarter (the \emph{cut}) and its rms over the fast part of the path,
    where the schedule's speed is above half its peak (the \emph{across-path error}). The \textbf{lag} is the time
    since the schedule passed the ring's point, positive late. The claim is scored on the frozen ring RMSE and the frozen
    tilt; the clock-free error and the lag are reported beside them.
    \item[Payload tilt] The angle between the ring's axis and the vertical: its mean over the hover and its
    peak after an event. The carry claim's peak tilt is over the whole flight, lift-off and landing included, and
    a peak over the path alone is reported beside it.
    \item[Disturbance response] For a constant force applied for a fixed time, the displacement along the
    force relative to its mean over the 5\,s before: steady (mean over the last 5\,s of the force), peak
    (while the force acts), overshoot (the largest excursion the other way after it ends) and time to
    recover (after which it stays within 1\,cm). \textbf{Stiffness:} the force over the steady displacement
    with no integral action.
% 11 Oct: the release metrics are the signed freeze (docs/freeze_2026-10-14.md section 1.1, Wesley 11 Oct 17:24,
% every question at its recommended option); the code is tools/release_bands.py. They replace the earlier "Detach"
% item (peak in 5 s, back in the hover's band, height change in 12 s).
    \item[Release metrics] For a UAV leaving the team, frozen on 11 October before the first flight they
    score; one rule for the twin and the rig. The ring's tilt is kept as a vector: with $\mathbf{n}$ the ring's
    body $z$ axis in the world frame, the tilt angle is $\arctan\!\big(\sqrt{n_x^2+n_y^2}\,/\,n_z\big)$ and
    the tilt vector $\boldsymbol{\theta}$ is that angle along the horizontal direction of $\mathbf{n}$, so that
    $|\boldsymbol{\theta}|$ is the scalar tilt above. The event time $t_e$ is the DETACH command for an announced
    release, the magnet's OFF for an unannounced one (the runner's RELEASE in the twin, the first magnet OFF in
    the rig's bag), and the first planner tick at which the leaver's planned tension falls for a ramped release;
    LAND is the LAND command. Over the pre-window $[t_e-5\,\mathrm{s},\,t_e)$ the mean
    $\bar{\boldsymbol{\theta}}_{\mathrm{pre}}$ and the spread
    $\sigma_{\mathrm{pre}}=\sqrt{\overline{|\boldsymbol{\theta}-\bar{\boldsymbol{\theta}}_{\mathrm{pre}}|^2}}$
    are taken. A row whose ring reference was not still for the 5\,s before the event (its height within
    5\,mm of its value at $t_e$ and its speed within 0.02\,m/s of its median over the last 2\,s) is flagged and
    still reported. \textbf{Tilt rise} $\Delta\theta$, the primary metric, is the maximum of
    $|\boldsymbol{\theta}(t)-\bar{\boldsymbol{\theta}}_{\mathrm{pre}}|$ over $[t_e,\,\mathrm{LAND}]$; when two
    runs are compared both are cut at the shorter hold, each 8 October pair at the rig flight's event-to-LAND,
    and the ledger's bands and the flights they score at $t_e+15$\,s. Beside it: the absolute peak,
    $\max|\boldsymbol{\theta}|$ over the same window, and the scalar rise, the absolute peak minus the mean of
    $|\boldsymbol{\theta}|$ over the pre-window. The 30$^\circ$ drop rule and the join bars stay on the
    absolute tilt. \textbf{Time back in band:} with $\bar{\boldsymbol{\theta}}_{\mathrm{post}}$ the mean over
    the last 5\,s before LAND, scored only when LAND comes at least 10\,s after $t_e$, and the half-width
    $b=\max(2\sigma_{\mathrm{pre}},\,0.5^\circ)$, it is the earliest $t$ at or after the time of $\Delta\theta$
    such that $|\boldsymbol{\theta}-\bar{\boldsymbol{\theta}}_{\mathrm{post}}|\le b$ for every sample in
    $[t,\,t+1\,\mathrm{s}]$, reported as $t-t_e$, or ``not back''. The dip (the largest drop of the ring's
    height below its pre-event mean), the detection time (from the magnet's OFF to the planner's resize), the
    steady offset, the landing lean, the clearances and the survivors' shares keep their scorers and are not
    part of the freeze.
    \item[Join metrics] For a newcomer admitted to the ring: the time from the last contact to the confirmation; the ring's
    tilt rise, on the release metric above with the contact as the event, from the contact to 8\,s after the
    confirmation; the ring dip over the same window; the hand-over step, the change in the newcomer's applied
    throttle, roll, pitch and yaw-rate commands between the 0.1\,s before the switch of authority and the 0.3\,s
    after it, as a fraction of its hover throttle and in degrees; the newcomer's tilt at the confirmation; and the
    clearances of the newcomer's body and rod to every carrier body and rod through its approach and, when it is
    taken back, through its take-back. For the full mission, also the tilt at the release, the mean tilt over the
    circle and the time from TAKEOFF to every UAV landed.
    \item[Magnet margin] For a release, the survivors' peak measured rod pull and the rate at which it rises,
    against the magnet's hold at that rod angle (Section~\ref{sec:hardware}); the one metric that says how close a
    transient came to letting a second magnet go.
    \item[Switch validation] For every resize, the planner's solver status and residual at the first live solve
    after the switch, the distance of that solve from an offline solve of the same logged problem run to
    convergence (not computed so far: the planner keeps a problem only when its solve fails), the solve time against the 60\,ms budget on the rig, and the step in each survivor's commanded
    throttle across the switch.
    \item[Solver health] The number of planner solves that failed or did not converge.
\end{description}
Each table cell is the mean over its runs, with the number of runs in brackets. Every arm of a lever table and
every ledger row also reports the state the ring was in before its event, the steady offset, the tilt and the
leaver's pull over the pre-window, so that a lever is read against the state it acted on and not only against
the other arm.

\paragraph{Resolving a lever pair.} A lever on the transient is tested as a one-variable pair of arms A and B,
each with at least two runs and one value per run ($\Delta\theta$ unless the table says otherwise), with the
sign of the difference written down before the runs. The pair is \emph{resolved} only if all three hold: the
sign of $\bar{B}-\bar{A}$ is the frozen sign; the arms' ranges do not overlap; and $|\bar{B}-\bar{A}|$ is at
least the larger of $1^\circ$ and twice the larger range. Otherwise it is \emph{unresolved}. A prediction of
``no sign'' is confirmed when the second or the third condition fails and falsified when both hold, in either
direction. Every row is labelled a \emph{replication} when both arms of a twin pair, or the rig flight of a
ledger row, were flown and registered before the rule was signed on 11 October, and a \emph{new prediction}
otherwise; the label is set at the signature and never changes.

% 11 Oct: the physical rationale for the bars (health check, item 3); the metric definitions above are the signed
% freeze of 11 Oct.
\paragraph{Why these bars.} Three physical limits say what a transient may do to the ring, and the bars of the
results sections sit inside them. The first is the magnets. A plate holds its magnet straight off the face by
about 9--10\,N and by less under shear, so a tilt that turns a rod's pull across its plate can let a magnet go:
every ring lost on the rig began with a magnet letting go with no command, two with the rods planned at
45$^\circ$ and two since the rods have been planned at 60$^\circ$ to keep each pull closer to its plate's axis
(Section~\ref{sec:res_leave}). [The hold against the angle of pull, still to be measured with a spring scale,
gives this limit its number.] The second is the fleet's own rule:
a ring tilt over 30$^\circ$, or a loss that leaves fewer than three carriers, drops the ring and lands every UAV
(Section~\ref{sec:safety}), so
30$^\circ$ is where the controller gives up, not a value to design to. The third is the model: the planner's
load--rod model and each tracker's rod pull assume a near-level ring, so a transient that stays within a few
degrees stays inside the regime the plan was computed for, and a transient that leaves it is being flown on a
plan that no longer describes the ring. The bars are therefore stated as a rise above the ring's own motion
before the event, small against the first two limits and inside the third, and every run is reported with its
measured rise so that a run that misses a bar is still read against the limits rather than against a number.
